\section{Relate Work}
\label{sec:Relatework}

\subsection{Graphic Design Generation}

Automatic graphic design generation has emerged as a crucial research area aimed at facilitating the design process.
CanvasVAE~\cite{yamaguchi2021canvasvae}  introduces a variational autoencoder to generate vector graphic designs represented as sets of canvas and element attributes. GOL~\cite{order} 
learns a design element ordering strategy to optimize the performance of design generative models. In recent years, a large body of methods on text-to-design generation has emerged to convert textual design intentions into multi-layer graphic designs, e.g., by building a cascaded system comprising multiple specialized models for different sub-tasks (e.g., background generation, typography generation)~\cite{cole,opencole},
formulating layered design generation as multi-layer transparent image generation~\cite{art}, 
utilizing a multimodal large language model to generate graphic designs as text-image documents~\cite{igd}, 
and training a simple diffusion model on sequences of discrete element attribute tokens ~\cite{LADEREIN}.There is also a line of work on automatic design composition, which aims to compose a set of input design elements into a complete design ~\cite{shabani2024visual,Graphist,LaDeCo}. 
FlexDM~\cite{flexdm} is a unified model that can solve multiple design tasks (e.g., image filling, layout generation, typography generation), by representing a graphic design as a set of elements, each comprising a set of fields, and formulating different design tasks as masked field prediction problems with different masking patterns.

Recent years have also seen a growing interest in developing models for generating layouts of graphic design elements~\cite{LayoutPrompter,jyothi2019layoutvae,Content-aware-Generative,VTN,PosterLlama,LayoutDiffusion,LayoutGAN,yang2026learninginteractionimagelayout}.
The success of Transformers in natural language processing has inspired their application to layout generation. LayoutTransformer~\cite{gupta2021layouttransformer} formulates layout generation as a sequence modeling problem and employs the self-attention mechanism to capture element relationships. BLT~\cite{kong2022blt} proposes a bidirectional Transformer that generates all layout tokens in parallel to enable flexible controllable generation.
Our method leverages Transformers to predict 3D human pose, 2D framing, and human image layout.

To our best knowledge, we are the first to model 3D human representations in 2D graphic designs, which has not been yet explored in previou studies.

\subsection{Human Pose Estimation and Generation}

3D Human pose and shape estimation
have been extensively studied in computer vision, with a range of applications in behavior analysis, animation and gaming.
Recent advances in expressive human pose and shape estimation (EHPS) enable the reconstruction of 3D whole-body meshes of a person\textemdash including the body, hands and face\textemdash from monocular images or videos using end-to-end models~\cite{OSX,cai2023smplerx,aios} 
that regress the parameters of a 3D parametric human model (e.g., SMPL-X~\cite{SMPLX}.

Conditional 3D human pose generation based on input conditions of various modalities has recently gained increasing attention. For example, ChatPose~\cite{feng2024chatpose} generates 3D human poses from natural language descriptions, optionally with image inputs. GenZI~\cite{li2024genzi} proposes a method to synthesize a posed 3D human that interacts with a given 3D scene according to an input text description. A recent study~\cite{jiang2024scaling} focuses on dynamic human-scene interaction modeling, generating 3D human motions (sequences of 3D poses) that are plausible in a given 3D scene context. NIFTY~\cite{kulkarni2024nifty} addresses the problem of object-conditioned human motion synthesis, which generates realistic 3D motions of a person interacting with a given 3D object.

In this work, we focus on generating 3D human poses in the context of graphic designs.

